\section{Introduction}\label{sec:intro}

Recency, frequency and monetary value (RFM) remain the most widely used basis for behavioural customer segmentation because they are cheap to compute, easy to explain and strongly predictive of future purchasing \citep{hughes1994,fader2005}. In practice each dimension is scored by quintile and customers are grouped either by score combinations or by clustering the scores or raw values \citep{chen2012,wedel2000}. \citet{rungruang2024} proposed an alternative: treat the fifteen quintile indicators (R1--R5, F1--F5, M1--M5) as a binary formal context and read segments off its concept lattice \citep{ganter1999}. Unlike clusters, formal concepts are defined by their descriptions (an intent such as $\{R5, F5, M5\}$), are ordered in a hierarchy, and come with association rules between them, which makes them directly usable by marketers.

The binary context has a cost that \citet{rungruang2024} themselves identify as future work: a customer just below a quintile cut point and one just above it receive disjoint attributes although their behaviour is almost the same. We find that this is not a corner case: 69--74\,\% of customers in three retail datasets lie within 0.05 of a quintile cut point on at least one dimension (Section~\ref{sec:res-interp}). For these customers the crisp segment depends on a near-arbitrary line, and small data changes move them between segments.

This paper replaces the binary context with a fuzzy one and builds segments from fuzzy formal concepts \citep{belohlavek2002,belohlavek2004}. Taken naively, the extension fails: under the Łukasiewicz residuum every customer belongs at least half-way to any concept whose intent contains only partial terms, so thousands of uninformative concepts pass the usual support threshold, and models built on them overfit (Section~\ref{sec:res-dev}). We make the extension work and test it thoroughly. Our contributions are:

\begin{enumerate}
\item \textbf{Method.} A fuzzy RFM context on a percentile-rank scale whose 0.5-cut is exactly the crisp quintile context; a \emph{core-support} threshold for fuzzy concepts; a greedy \emph{summary} that selects a small set of concepts; and a \emph{tie-aware} membership for heavily tied frequencies.
\item \textbf{Theory.} Crisp RFM-FCA is a special case of the method; core support reduces to ordinary support on binary data and is anti-monotone; the Łukasiewicz floor explains why sigma-count support fails; and the summary objective is monotone submodular, so the greedy selection is within $1-1/e$ of optimal \citep{nemhauser1978}.
\item \textbf{Evidence.} A replication of \citet{rungruang2024} and a rolling-origin, out-of-sample comparison on three datasets against the crisp method, fuzzy memberships alone, log-RFM, and K-means, Gaussian-mixture and fuzzy C-means clustering, covering prediction, targeting value, robustness, stability, structural interpretability and sensitivity to every parameter fixed early. Analysis plans were committed to version control before each analysis, and the two later datasets serve as confirmation.
\end{enumerate}

The results are positive but bounded. Fuzzy RFM-FCA beats crisp RFM-FCA on every dataset and in 47 of 48 sensitivity settings, reaches more revenue in the most valuable customer decile and gives the most robust and stable segments of all methods. It is not a better predictor than a plain log-RFM model, it is roughly level with fuzzy C-means, and its full lattice has about four times as many segments as the crisp one, which the 10-concept summary addresses. We report these limits alongside the gains.

Section~\ref{sec:related} reviews related work, Section~\ref{sec:method} presents the method and its properties, Section~\ref{sec:design} the experimental design, Section~\ref{sec:results} the results, and Section~\ref{sec:discussion} their implications and limitations.
